\chapter{The Direction}
% OUTLINE Ch3 "The Direction (linear algebra at a point)". Laws: CITATION
% (Grassmann 1844; Ricci & Levi-Civita 1901; Dirac 1930; Lee/standard for
% the modern formalization), LENS (zero device appearances), Vol-1 laws.
% The two transformation laws DERIVED from the chain rule, never decreed.
% No metric: vectors and covectors stay DISTINCT species; raising/lowering
% flagged as Ch7's purchase. NO series jargon (tange/funge stays in Vol 1;
% the reader's own words are the raised and lowered indices). Gate: Kodo.

Stand at one point of the last chapter's manifold and ask the physicist's
most pointed question: a body passes through here, this instant --- which
way, and how fast? The answer exists before any observer lays down
coordinates. The thrown stone does not consult a chart to have a velocity;
the velocity is a fact at the point, an arrow with a direction and a rate,
and every chart that covers the point merely reports it in that chart's own
components. This chapter is about the arrow, the read-outs that measure the
arrow, and the discovery --- the physicist has been using it in index
notation for a century --- that arrows and read-outs respond to a change of
observer in two opposite ways, and must never be confused.

First the arrow, built honestly from what previous chapters own. A curve
through the point --- a smooth path passing through it at time zero --- has,
in any chart, a column of component-rates: differentiate the coordinates
along the path, the classical calculus of number-space, which the treaty of
Chapter 2 licenses region by region. Two curves through the point may have
the same component-rates in one chart; the transition maps being smooth,
they then have the same component-rates in every chart --- agreement is
chart-independent, and so it is a fact about the point, not the furniture.
A velocity is exactly such an agreement-class of curves, and the velocities
at a point form a vector space: they add, they scale, and the dimension is
the chart's dimension. The trade calls it the tangent space at the point;
the modern formalization is standard and lives in the same references as
the last chapter's (Lee's, for this series' purposes). The idea that
directed magnitudes in $n$ dimensions deserve their own algebra --- addition
of arrows, scaling of arrows, without any appeal to lengths --- is far
older, and its author paid for being early: Hermann Grassmann published the
calculus of extension in 1844, a schoolmaster writing outside the
universities, and the mathematics of his century mostly declined to read
him. This series cites him the way the trade eventually did: as the one who
had the vector space first.

Now the read-out. An arrow is only half of a measurement; the other half is
the thing that turns an arrow into a number. Let a quantity be defined near
the point --- a temperature, a potential, any scalar field the bench
provides --- and ask: at what rate does a body with this velocity see that
quantity change? The answer is a number for each velocity, and it depends
on the velocity linearly: twice the speed, twice the rate; the sum of two
motions, the sum of the rates. A linear machine from arrows to numbers is
called a covector, and the covectors at a point form their own vector
space, of the same dimension, called the dual. The physical reading should
be said plainly: covectors are the measuring devices of this subject. An
arrow is a state of motion; a covector is an instrument reading applied to
motion; and the pairing of the two --- instrument applied to state, number
out --- is the elementary act from which everything later in this book is
assembled. The physicist has met the structure in another guise: Dirac's
formulation of the quantum theory, in his 1930 \emph{Principles}, pairs
kets with bras --- states with the linear functionals that read them ---
and the bracket that gives the formalism its name is exactly this pairing,
the dual space working in the physicist's own house.

Here is the chapter's centerpiece, and it is a derivation, not a decree.
Change observers: a second chart covers the point, and the treaty supplies
the dictionary --- the transition map --- with its matrix of partial
derivatives, the Jacobian of the change, evaluated at the point. Follow
the arrow first. Its new components come from differentiating the new
coordinates along the same curve, and the chain rule --- the same chain
rule the reader has always used --- says: new components equal the Jacobian
applied to old components. Now follow the read-out. Its job is to report
the same physical rate no matter the observer, so its components must
change in exactly the compensating way: by the inverse of the Jacobian,
arranged so that when instrument meets state the Jacobians cancel and the
number --- the measured rate --- comes out unchanged. Two laws, then, and
nobody chose them: components of arrows transform one way, components of
read-outs the opposite way, and both fall out of the chain rule applied to
the observers' dictionary. The trade's names are contravariant for the
arrow's law and covariant for the read-out's, and the physicist's raised
and lowered indices are precisely the bookkeeping of which law a given
column of components obeys. The calculus built on exactly this distinction
--- systems of components classified by their transformation law, with
rules for combining them so that measured numbers stay observer-invariant
--- was published by Ricci and Levi-Civita in 1901, the absolute
differential calculus, and it is the machinery general relativity would
later be written in. Nothing about it was invented for relativity; it was
waiting.

One more honest note completes the chapter, and it continues this book's
longest-running theme. The reader's working habit converts freely between
raised and lowered indices --- between arrows and read-outs --- and nothing
in this chapter licenses the conversion. As matters stand, the two are
different species: an arrow is a state, a read-out is an instrument, their
spaces are merely of equal dimension, and no natural dictionary matches
them element to element. The conversion the physicist performs daily is a
purchase, and the coin it is bought with is precisely the ruler this
construction has not yet admitted: a metric. When the metric arrives,
chapters from now, it will do exactly this work --- pair each arrow with
its own read-out --- and the reader will see index-lowering for what it is:
not notation, but geometry, at a price the book will pay in the open. Until
then the species stay separate, and the separation is not pedantry; it is
the reason the two transformation laws exist at all.

What comes next is multiplication. One arrow in, one number out, was the
covector; the objects of physics --- the stresses, the fields, the
curvatures this series is walking toward --- take several arrows and
read-outs at once, in slots, and answer with a number, linearly in each
slot. The machines with slots are the tensors, the subject of the next
chapter, and their first and oldest example was not invented by a geometer
at all: it was found inside a loaded body, by Cauchy, holding the stresses
of continuum matter --- which is where the fourth station of this series'
preface will briefly come back into view.
